\chapter{Series and Finite Remainders}\label{ch:infinite-series}

An infinite sum is useful here only when a finite prefix comes with a
bound on every finite future correction.  This produces a number
computation by stabilization.  Algebraic identities between prefixes,
combined with tail bounds, then prove identities between the computations.

\section{Geometric and alternating sums}
\label{thm:finite-geometric-series-interface}
For rational $|r|<1$,
\[
 \left|\sum_{k=N+1}^{M}r^k\right|
       \le\frac{|r|^{N+1}}{1-|r|}\qquad(M>N).
\]
The right side is a rational error tending to zero.  Thus the geometric
series computes $1/(1-r)$ by its finite identity.

If $a_n$ are nonnegative, decrease to zero with an effective bound, then
even and odd partial sums of $\sum(-1)^na_n$ interleave; pairing adjacent
terms proves that the next term bounds the remainder.  Computed rather
than rational terms are evaluated to a finite total error budget.

\begin{example}[Divergence is also a computation]
\label{thm:harmonic-dyadic-lower-bound}
Grouping the terms with $2^{j-1}<n\le2^j$ gives
\[
 \sum_{n=1}^{2^m}\frac1n\ge1+\frac m2.
\]
For every rational height, this produces an explicit finite partial sum
above it.  No infinite value is introduced.
\end{example}

The sum of reciprocal squares has the independent tail bound
\[
 \sum_{n=N+1}^{M}\frac1{n^2}
 \le\sum_{n=N+1}^{M}\left(\frac1{n-1}-\frac1n\right)\le\frac1N.
\]
This constructs the number before identifying it with $\pi^2/6$.
The identification will be a Fourier computation in
Chapter~\ref{ch:fourier}.

\section{Arctangent: a finite identity integrates to a series}
\label{thm:leibniz-arctan}
The exact polynomial identity
\[
 \frac1{1+t^2}=\sum_{k=0}^{N-1}(-1)^kt^{2k}
              +\frac{(-1)^Nt^{2N}}{1+t^2}
\]
has an integrable remainder.  On $[0,x]$ with $0\le x\le1$ its
absolute integral is at most $x^{2N+1}/(2N+1)$.  Finite polynomial
integration and the already proved angle-area comparison yield
\[
 A(x)=\sum_{k=0}^{N-1}\frac{(-1)^kx^{2k+1}}{2k+1}+R_N(x),
 \qquad |R_N(x)|\le\frac{x^{2N+1}}{2N+1}.
\]
At $x=1$ this is the Leibniz computation of $\pi/4$.  This is a proof of
agreement between a geometric computation and a series computation, not
an identification by notation.

\section{A power series is a prefix together with a majorant}
Suppose the coefficients $a_n$ are computable, and rational constants
$M,R>0$ give $|a_n|\le M R^{-n}$.  On $|x|\le r<R$, put $q=r/R$.
Then
\[
 F_N(x)=\sum_{n=0}^Na_nx^n,\qquad
 |F-F_N|\le\frac{M q^{N+1}}{1-q}.
\]
The inequality means a bound on all finite future tails, and so constructs
$F$.  Coefficient evaluations are made accurate enough that their finite
weighted error is smaller than a prescribed remaining budget.  The same
construction works in rational complex boxes on a smaller disk.

For the differentiated prefixes the tail is bounded by
\[
 \frac{M}{R}q^N\left(\frac{N+1}{1-q}+\frac{q}{(1-q)^2}\right).
\]
This follows by summing $nq^{n-1}$, or differentiating a finite geometric
identity and bounding its final terms.  Higher differentiated tails follow
in exactly the same finite way, with explicit polynomial factors in $n$.
One fixes the required order and writes its bound; no unrestricted exchange
of limits is presumed.

\begin{theorem}[Differentiation and integration of a supplied series]
\label{thm:finite-taylor-ftc-prefix}
On a smaller interval $[-r,r]$, the derivative is the computation with
coefficients $(n+1)a_{n+1}$.  Its integral is computed term by term, and
\[
 \int_a^b\sum_{n=0}^{\infty}a_nx^n\,dx
   =\sum_{n=0}^{\infty}\frac{a_n}{n+1}(b^{n+1}-a^{n+1}).
\]
\end{theorem}
\begin{proof}
Each finite polynomial has its exact algebraic derivative and primitive.
For all prefixes the second derivatives have the common bound
\[
 K=\frac{2M}{R^2(1-q)^3}.
\]
The finite polynomial remainder identity gives
$|F_N(y)-F_N(x)-F_N'(x)(y-x)|\le K|y-x|^2/2$ for a segment in
$[-r,r]$.  One can derive this identity by integrating the polynomial
second derivative twice; polynomial integration is just the binomial
identity and finite telescoping.  Pass to the computations using the
explicit tails for $F_N$ and $F_N'$.  This proves the derivative assertion.
Uniform tail bounds change an integral by at most its interval length
times the tail, proving the second assertion.
\end{proof}

For real coefficients, $F_N+Kx^2/2$ is convex, by the same finite polynomial
second-difference identity.  Passing its inequalities through the tails
proves that $F+Kx^2/2$ is convex.  Thus this derivative agrees with the
convex-secant computation of Chapter~\ref{ch:effective-calculus}; it is not
a competing notion of derivative.

\section{Products, coefficient identities, and Taylor errors}
For two supplied series, multiply finite prefixes.  Absolute majorants
bound the two omitted tails and their product.  Regrouping their finite
terms proves the Cauchy product identity.  The same argument in two
variables justifies the coefficient comparisons used for exponential
addition below.

A Taylor polynomial is another finite computation of the same function.
When a translated series or a finite remainder bound has been supplied,
its error is bounded by that data.  For example the uniform bound $K$
above gives the first-order error $K h^2/2$.  More generally a supplied
bound $B$ on the polynomial/series $(m+1)$st derivatives gives
$B|h|^{m+1}/(m+1)!$ by repeated finite polynomial integration and its
uniform tails.  We do not infer such a remainder theorem from pointwise
rational derivatives alone.

\section{Binomial square-root computations}
Put
\[
 c_0=1,\qquad c_{n+1}=\frac{2n+1}{2n+2}c_n
       =\frac{\binom{2n+2}{n+1}}{4^{n+1}}.
\]
The coefficients are rational and $0<c_n\le1$, so
$B(z)=\sum_{n\ge0}c_nz^n$ has a geometric tail on $|z|\le r<1$.
The recurrence gives the coefficient identity
\[
 2(1-z)B'(z)=B(z).
\]
By the product rules for series, every nonconstant coefficient of
$(1-z)B(z)^2$ has derivative zero, and its constant coefficient is $1$.
Consequently
\[
 (1-z)B(z)^2=1.
\]
For $0\le z<1$, positivity of the coefficients identifies this computation
with the positive inverse square root $(1-z)^{-1/2}$.  The argument is
coefficient algebra with tails; it does not use an ODE uniqueness theorem
for arbitrary rational-point solutions.

This example will reappear in the construction of an elliptic integral.
It also illustrates how a recurrence can supply a special function, while
an algebraic identity identifies what the computation evaluates.
